%% !TEX root = ../manual_gaolv5.tex
% Chapter 7.
%% Chapter 7: Interval arithmetic: the operations. Sources: gaol.tex of f5bf180
%% lines 3493-4445 (chapter "Interval Arithmetic": the introduction and the
%% domains, functional arithmetic, relational arithmetic, cancellative addition
%% and subtraction), doc/using.md lines 695-842 (tools of interval algorithms,
%% the reverse functions of pow, max, min, sign and floor).
\chapter{Interval arithmetic: the operations}\label{chap:arithmetic}

An interval program is written with the operators and the functions of C++,
which GAOL overloads for the type \code{interval}: a sum, a product, a square
root or a cosine of intervals is an interval that contains every value the
operation takes on the real numbers of its arguments. This chapter is the
reference of these operations. The containment property they satisfy, the
roundings that give their bounds and the mathematical libraries that compute
the elementary functions are the subject of Chapter~\ref{chap:concepts}. Here
we state, for each operation, what it computes, on which domain, which
operation of IEEE 1788-2015 \cite{IEEE1788:2015} it is, and what differs from
GAOL~4.

\medskip
The operations are of two kinds. A \emph{functional} operation extends a
function of real numbers to intervals: the square root of an interval is the
hull of the square roots of its points. A \emph{relational} operation inverts
a relation: the relational square root of $J$ with respect to $I$ is the hull
of the $x$ of $I$ whose square is in $J$, the negative ones included. The
relational operations are the reverse functions of IEEE 1788-2015, and the
tools with which the algorithms of constraint programming contract their
domains.

\medskip
The first section recalls the containment principle, the two kinds of
operations and the rule of the domains. The second section is the reference
of the functional arithmetic: the basic operations, the powers and roots, the
exponentials and logarithms, the trigonometric functions, the hyperbolic
functions and the other functions of the standard. The third section is the
reference of the relational arithmetic, and gives its correspondence with the
reverse functions of the standard. The fourth section describes the
cancellative addition and subtraction. The tightness of each result is
stated in Chapter~\ref{chap:accuracy} (\pxref{sec:accuracy-table}), and the
names of the standard are the subject of Chapter~\ref{chap:ieee1788}.

\section{Functional and relational operations}\label{sec:arithmetic-kinds}
\cindex{functional arithmetic}\cindex{relational arithmetic}%

The containment principle of floating-point interval arithmetic
(\pxref{sec:containment}) requires that, for any operation $\circ$ and any
intervals $I$ and $J$,
\begin{equation*}
I\circ J=\hull{\{i\circ j\mid i\in I, j\in J\}}
\end{equation*}
where \hull{} maps any set of real numbers to the smallest floating-point
interval that contains it. The interval square root, for instance, is
\begin{equation*}
\sqrt{I}=\hull{\{\sqrt{i}\mid i\in I\}}
\end{equation*}
and, the square root being increasing, $\sqrt{\itv{1}{2}}=\itv{1}{\sqrt{2}}$.
This is the \dfn{functional square root}: it extends the function
$x\mapsto\sqrt{x}$ to intervals. Another reading of the square root gives
both the negative and the positive values.

\begin{Def}\label{def:arithmetic-relational-sqrt}
The \dfn{relational square root} of an interval $I$ is the hull of the real
numbers whose square is in $I$:
\begin{equation*}
\sqrt{_rI}=\hull{\{j\in\RSet\mid\exists i\in I\colon j^2=i\}}.
\end{equation*}
\end{Def}

\noindent Hence $\sqrt{_r\itv{1}{2}}=\itv{-\sqrt{2}}{\sqrt{2}}$.

\begin{Ex}
The relational square root arises from the relation $x^2=y$ between two
variables, which a program may want to write $x=\sqrt{_ry}$ to compute $x$
from $y$. The functional square root is not suitable there: its result lies
in \itv{0}{+\infty}, and the negative values of $x$ would be lost.
\end{Ex}

Some applications, mainly in constraint programming, need such operators.
GAOL therefore offers both the functional and the relational versions of the
main arithmetic operators (\pxref{sec:relational-idea}); a relational
operator contracts one operand of a relation from the others, which is how
the interval algorithms of Chapter~\ref{chap:overview} use it
(\pxref{sec:tools}).

\paragraph{Domains.}
A function is computed on the part of its argument that lies within its
domain, which is the domain of IEEE 1788-2015 (Tables~9.1 and 10.5): the
part outside is left out, and an argument with no point of the domain gives
the empty set, as \code{log(\itv{-4}{0})} does.\newinvfive{} Every function
gives the empty set for an empty argument. How tight each result is, for each
operation, is stated in Chapter~\ref{chap:accuracy}
(\pxref{sec:accuracy-table}).
\cindex{domain of a function}%

\section{Functional arithmetic}\label{sec:arithmetic-functional}

Each operation of this section is the hull of the image, by the function of
the same name, of the part of its arguments within its domain. The entries
give, for each one, the operation of IEEE 1788-2015 it implements, the
domain where it is not the whole real line, and what GAOL~v5 changed.

\subsection{Basic operations}\label{sec:arithmetic-basic}

\defmethodx{interval}{interval}{operator+}{(void) const}
\deffunx{interval}{operator+}{(const interval\& \var{I}, double \var{d})}
\deffunx{interval}{operator+}{(double \var{d}, const interval\& \var{I})}
\begin{deffun}{interval}{operator+}{(const interval\& \var{I1}, const interval\& \var{I2})}
  Addition of two intervals, or of an interval and a double: \ieee{add} of
  IEEE 1788-2015.
\end{deffun}

\defmethodx{interval}{interval}{operator-}{(void) const}
\deffunx{interval}{operator-}{(const interval\& \var{I}, double \var{d})}
\deffunx{interval}{operator-}{(double \var{d}, const interval\& \var{I})}
\begin{deffun}{interval}{operator-}{(const interval\& \var{I1}, const interval\& \var{I2})}
  Negation (\ieee{neg}), or subtraction (\ieee{sub}) of two intervals, or of
  an interval and a double.
\end{deffun}

\deffunx{interval}{operator*}{(const interval\& \var{I}, double \var{d})}
\deffunx{interval}{operator*}{(double \var{d}, const interval\& \var{I})}
\begin{deffun}{interval}{operator*}{(const interval\& \var{I1}, const interval\& \var{I2})}
  Multiplication of two intervals, or of an interval and a double: \ieee{mul}
  of IEEE 1788-2015. A zero bound times an infinite bound counts as $0$, as
  the product of a real number of the interval by 0.\newinvfive{} \code{d * I}
  is computed as \code{I * d}, and \code{n * I} of an integer \code{n} as
  \code{I * n}: the bounds are the same, with the same signs of a zero bound
  (\pxref{sec:zero-sign}).\newinvfive
\end{deffun}

\deffunx{interval}{operator/}{(const interval\& \var{I}, double \var{d})}
\deffunx{interval}{operator/}{(double \var{d}, const interval\& \var{I})}
\begin{deffun}{interval}{operator/}{(const interval\& \var{I1}, const interval\& \var{I2})}
  Functional division of two intervals, or of an interval and a double:
  \ieee{div} of IEEE 1788-2015. The quotient by an interval containing 0 is
  the hull of the quotients by its nonzero part, and the quotient by
  \itv{0}{0} is the empty set.
\begin{example}
cout << interval(1,2) * interval(-3,4) << ' '
     << interval(0,1) * interval(1,GAOL_INFINITY) << ' '
     << interval(1,2) / interval(0,1) << ' '
     << interval(1,2) / interval(-1,1) << ' '
     << interval(1,2) / interval::zero();
@outputs^[-6, 8] [0, inf] [1, inf] [-inf, inf] [empty]~
\end{example}
\end{deffun}

\deffunx{interval}{operator\%}{(const interval\& \var{I}, double \var{d})}
\deffunx{interval}{operator\%}{(double \var{d}, const interval\& \var{I})}
\begin{deffun}{interval}{operator\%}{(const interval\& \var{I1}, const interval\& \var{I2})}
  Relational division of two intervals, or of an interval and a double
  (\pxref{sec:relational-arithmetic}).
\end{deffun}

\defmethodx{interval}{interval}{inverse}{(void) const}
\begin{deffun}{interval}{inverse}{(const interval\& \var{I})}
  \begin{operation}
    \op{$1/I$}
  \end{operation}
  Inverse of \var{I} (of \code{self} for the method): \ieee{recip} of
  IEEE 1788-2015. The inverse of an interval containing 0 is the hull of the
  inverses of its nonzero part.
\begin{example}
cout << inverse(interval(2,4)) << ' ' << inverse(interval(-1,2));
@outputs^[0.25, 0.5] [-inf, inf]~
\end{example}
\end{deffun}

\begin{deffun}{interval}{sqr}{(const interval\& \var{I})}
  \begin{operation}
    \op{$\var{I}^2$}
  \end{operation}
  Square of \var{I}: \ieee{sqr} of IEEE 1788-2015.
\end{deffun}

\begin{deffun}{interval}{sqrt}{(const interval\& \var{I})}
  \begin{operation}
    \op{$\sqrt{I}$}
  \end{operation}
  Functional square root of \var{I}, on the part of \var{I} in
  \itv{0}{+\infty}: \ieee{sqrt} of IEEE 1788-2015. The bounds are the
  tightest ones, whatever the rounding of the square root of the C
  library.\newinvfive
\begin{example}
cout << sqrt(interval(-4,9)) << ' ' << sqrt(interval(-4,-1));
@outputs^[0, 3] [empty]~
\end{example}
\end{deffun}

\begin{deffun}{interval}{fma}{(const interval\& \var{X}, const interval\& \var{Y}, const interval\& \var{Z})}
  \begin{operation}
    \op{$\hull{\{xy+z\mid x\in X, y\in Y, z\in Z\}}$}
  \end{operation}
  Fused multiply-add: \ieee{fma} of IEEE 1788-2015. Each bound is rounded
  once, and the result is the tightest enclosure.\newinvfive
\begin{example}
cout << fma(interval(1,2), interval(3,4), interval(-1,1));
@outputs^[2, 9]~
\end{example}
\end{deffun}

\subsection{Powers and roots}\label{sec:powers}
\cindex{power}\cindex{root}%

The functions of this subsection are those of the namespace \code{gaol}. The
namespace \code{gaol\_ieee1788} has its own \code{pow}, the \ieee{pow} of
the standard: the two powers are compared in Chapter~\ref{chap:overview}
(\pxref{sec:pow-namespaces}).

\deffunx{interval}{pow}{(const interval\& \var{I}, int \var{n})}
\deffunx{interval}{pow}{(const interval\& \var{I}, unsigned int \var{n})}
\deffunx{interval}{pow}{(const interval\& \var{I}, double \var{p})}
\deffunx{interval}{pow}{(const interval\& \var{I}, long \var{n})}
\deffunx{interval}{pow}{(const interval\& \var{I}, long long \var{n})}
\begin{deffun}{interval}{pow}{(const interval\& \var{I}, const interval\& \var{J})}
  \begin{operation}
    \op{$I^n$}
    \op{$I^p$}
    \op{$I^J$}
  \end{operation}
  Power function. \code{pow} is hybrid in GAOL v5.\newinvfive
  \begin{itemize}
  \item[$\bullet$] \textbf{An integer exponent} takes the \emph{integer
    power}, \ieee{pown} of IEEE 1788-2015: \var{I} to the power \var{n},
    the negative part of \var{I} included, where $x^0=1$ for every $x$, 0
    included, and $x^n$ has no value at 0 for $n<0$. A double \var{p} and a
    degenerate interval $\var{J}=\itv{n}{n}$ whose value is an integer take
    the integer power as well, beyond the \code{int}s too.\newinvfive{} An
    exponent \var{n} of another integer type (\code{long}, \code{long long},
    \code{short}\dots), with which \code{pow} was ambiguous, gives the same
    integer power; beyond 64 bits (an \code{\_\_int128}), it gives the power of $\pm(2^{63}-1)$ or
    $\pm(2^{63}-2)$, with the sign and the parity of \var{n}, every power of
    a double but $0$, $\pm1$ and $\pm\infty$ being beyond the doubles
    there.\newinvfive
  \item[$\bullet$] \textbf{Any other exponent} takes the \emph{real power},
    \ieee{pow} of IEEE 1788-2015, on the part of \var{I} in \itv{0}{+\infty},
    where $0^y$ is $0$ for $y>0$ and has no value for $y\leq0$. An exponent
    \itv{+\infty}{+\infty} or \itv{-\infty}{-\infty}, or an infinite or NaN
    \var{p}, gives the empty set.
  \end{itemize}
  The integer powers are computed from exact products within the
  \code{int}s, $x^{-n}$ included. They are the tightest enclosures, or one
  double beyond where the power is closer to a double than twice a bound of
  their error proved beforehand: $2^{-98}$ for $n=3$, $2^{-88}$ for
  $n=1000$, $2^{-62}$ at most, relatively. They are exact where the power is
  a double. Beyond the \code{int}s, the integer powers come from the power of
  CORE-MATH at the bounds of \var{I}: they are the tightest enclosures where
  $n$ is a double, and within three doubles of them
  otherwise.\newinvfive{} The real power takes the power of CORE-MATH at the
  corners of the box $I\times J$ where it is the least and the greatest. The
  upper bound is the tightest one, and so is the lower bound: it is the power
  itself where that power is a double, as $4^{0.5}$, $9^{1.5}$ or $2^{-3}$,
  and the double below the value of CORE-MATH otherwise. At a corner where
  the base is $0$ or $+\infty$ or the exponent is infinite, the bound is the
  limit of $x^y$ there, $0$ or $+\infty$ ($1$ for a base $1$ or an exponent
  $0$): \code{pow(\itv{4}{+\infty}, 0.5)} is $[2, +\infty]$.\newinvfive
\begin{example}
cout << pow(interval(-4,-1), 2) << ' '
     << pow(interval(-4,-1), interval(2.0)) << ' '
     << pow(interval(-4,9), interval(0.5)) << ' '
     << pow(interval(-4,-1), interval(1.5, 2.5)) << endl;
@outputs^[1, 16] [1, 16] [0, 3] [empty]~
cout << pow(interval(0.0), interval(0.0)) << ' '
     << pow(interval(-2,2), -2) << ' '
     << pow(interval(4), 0.5) << endl;
@outputs^[1] [0.25, inf] [2]~
\end{example}
\end{deffun}

\begin{Rem}\label{rem:arithmetic-pow-monotone}
The result of \code{pow} is not monotone for the inclusion of \var{J}:
\code{pow(\itv{-4}{-1}, \itv{2}{2})} is \itv{1}{16}, and
\code{pow(\itv{-4}{-1}, \itv{1.5}{2.5})} the empty set. The result for a
degenerate integer exponent $[n]$ is not within the result for an exponent
around $n$, and \code{pow} is not one function of $x$ and $y$. GAOL~4
truncated a double exponent \var{p}, and computed the powers of a negative
base on its magnitude.
\end{Rem}

\deffunx{interval}{nth\_root}{(const interval\& \var{I}, unsigned int \var{n})}
\deffunx{interval}{nth\_root}{(const interval\& \var{I}, long long \var{q})}
\deffunx{interval}{nth\_root}{(const interval\& \var{I}, unsigned long long \var{q})}
\begin{deffun}{interval}{nth\_root}{(const interval\& \var{I}, int \var{q})}
  \begin{operation}
    \op{$^{n}\sqrt{I}$}
  \end{operation}
  Computes the \var{n}th functional root of \var{I}: \ieee{rootn} of
  IEEE 1788-2015. For an odd \var{n}, the root of a negative number is the
  opposite of the root of its magnitude, and the root is defined on the whole
  real line; for an even \var{n}, only the part of \var{I} in
  \itv{0}{+\infty} is considered.\newinvfive{} A negative \var{q} gives
  $1/x^{1/|q|}$, defined on $\RSet\setminus\{0\}$ for an odd \var{q} and on
  $\left]0,+\infty\right]$ for an even one. The function returns the empty
  set for $n=0$.

  The bounds are the tightest ones, or up to 2 doubles beyond (4 for a
  negative \var{q}): each is proved with integer powers. The cube root is
  that of CORE-MATH, the tightest one, and the root of a double that is the
  \var{n}th power of a double is that double. An order of another integer
  type (\code{long}, \code{long long}\dots), with which \code{nth\_root} was
  ambiguous, is taken as an \code{int} or an \code{unsigned int} where it is
  one. Beyond them, the roots are enclosed by \code{pow} with the exponent
  $1/q$ taken as an interval, $x^{1/q}$ being within $2^{-21}$ of 1 for every
  double $x>0$ there, after one check of the rounding direction for a
  \code{long long} or an \code{unsigned long long} order.\newinvfive
\begin{example}
cout << nth_root(interval(-8,27), 3) << ' '
     << nth_root(interval(-4,16), 2) << ' '
     << nth_root(interval(4,16), -2) << ' '
     << nth_root(interval(-8,-1), -3);
@outputs^[-2, 3] [0, 4] [0.25, 0.5] [-1, -0.5]~
\end{example}
\end{deffun}

\subsection{Exponentials and logarithms}\label{sec:arithmetic-exp-log}
\cindex{exponential}\cindex{logarithm}%

The functions of this subsection, and of the next two, are bounded with
CORE-MATH at the bounds of the interval, or at the points within it where
the function reaches its extrema. CORE-MATH being correctly rounded, the
value at a bound, computed rounding upward, is the upper bound, and the
double below it is the lower bound. These are the tightest bounds, unless the
exact value is a double; the functions give that double where they know that
it happens ($\exp(0)=1$, $\log(1)=0$, $\sin(0)=0$, $2^3=8$\ldots).\newinvfive
\cindex{CORE-MATH}%

\begin{deffun}{interval}{exp}{(const interval\& \var{I})}
  Exponential of \var{I}.
\end{deffun}

\deffunx{interval}{exp2}{(const interval\& \var{I})}
\begin{deffun}{interval}{exp10}{(const interval\& \var{I})}
  $2^x$ and $10^x$ for $x$ in \var{I}: \ieee{exp2} and \ieee{exp10} of IEEE
  1788-2015.\newinvfive
\end{deffun}

\deffunx{interval}{expm1}{(const interval\& \var{I})}
\deffunx{interval}{exp2m1}{(const interval\& \var{I})}
\begin{deffun}{interval}{exp10m1}{(const interval\& \var{I})}
  $e^x-1$, $2^x-1$ and $10^x-1$ for $x$ in \var{I}, within
  \itv{-1}{+\infty}: \ieee{expm1}, \ieee{exp2m1} and \ieee{exp10m1} of IEEE
  1788-2015 (Table~10.5).\newinvfive
\end{deffun}

\begin{deffun}{interval}{log}{(const interval\& \var{I})}
  Natural logarithm of \var{I}, defined on $\left]0,+\infty\right]$. An
  interval holding no positive number gives the empty set, where GAOL~4 kept
  its part in \itv{0}{+\infty} and gave \itv{-\infty}{-\text{MAX}} for
  \code{log(\itv{-4}{0})}.\newinvfive
\end{deffun}

\deffunx{interval}{log2}{(const interval\& \var{I})}
\begin{deffun}{interval}{log10}{(const interval\& \var{I})}
  Logarithms in base 2 and 10 of \var{I}: \ieee{log2} and \ieee{log10} of
  IEEE 1788-2015.\newinvfive
\end{deffun}

\deffunx{interval}{log1p}{(const interval\& \var{I})}
\deffunx{interval}{log2p1}{(const interval\& \var{I})}
\begin{deffun}{interval}{log10p1}{(const interval\& \var{I})}
  $\log(1+x)$, $\log_2(1+x)$ and $\log_{10}(1+x)$ for $x$ in \var{I}, defined
  on $\left]-1,+\infty\right]$, $-\infty$ being their limit at $-1$:
  \ieee{logp1}, \ieee{log2p1} and \ieee{log10p1} of IEEE 1788-2015
  (Table~10.5).\newinvfive
\begin{example}
cout << exp(interval(0.0)) << ' ' << log(interval(1)) << ' '
     << log(interval(-4,0.0)) << ' '
     << log(interval(-4,1)) << endl;
@outputs^[1] [0] [empty] [-inf, 0]~
cout << exp2(interval(3)) << ' ' << log2(interval(0.25)) << ' '
     << log10(interval(0.0,1000)) << ' '
     << expm1(interval(-1e300,0.0)) << ' '
     << log1p(interval(-2,0.0)) << endl;
@outputs^[8] [-2] [-inf, 3] [-1, 0] [-inf, 0]~
\end{example}
\end{deffun}

\subsection{Trigonometric functions}\label{sec:arithmetic-trigonometric}
\cindex{trigonometric functions}%

\begin{deffun}{interval}{cos}{(const interval\& \var{I})}
  Returns the cosine of \var{I}, within \itv{-1}{1}.
\end{deffun}

\begin{deffun}{interval}{sin}{(const interval\& \var{I})}
  Returns the sine of \var{I}, within \itv{-1}{1}.
\end{deffun}

\begin{deffun}{interval}{tan}{(const interval\& \var{I})}
  Returns the tangent of \var{I}: \itv{-\infty}{+\infty} when \var{I} holds a
  pole. No double is a pole: \code{M\_PI\_2}, the double below $\pi/2$, is
  $6.1\cdot10^{-17}$ short of it. The interval
  \itv{-\code{M\_PI\_2}}{\code{M\_PI\_2}} holds no pole and has a finite
  tangent (the fourth result of the example); its width is the double below
  $\pi$, and the intervals of width between that double and $\pi$ are a
  theoretical case.

  The three functions find the pieces of \var{I} on which they are monotonic
  by dividing its bounds by an enclosure of $\pi$. Where the quotients cannot
  tell, next to an extremum or a pole and at the large magnitudes, they use
  the signs of the derivative at the bounds, which CORE-MATH gives exactly.
  Their bounds are within one double of the tightest ones at every
  magnitude.\newinvfive{} GAOL~4 gave \itv{-1}{1} for
  \code{cos(\itv{2^{60}}{2^{60}})}, and computed the sine as a shifted
  cosine, which widened narrow intervals.
\begin{example}
cout << cos(interval(1,4)) << endl
     << sin(interval(1,2)) << endl
     << tan(interval(1,2)) << endl
     @rem^// 1.5707963267948966 is M_PI_2, the double below pi/2~
     << tan(interval(-1.5707963267948966, 1.5707963267948966)) << endl
     @rem^// 1152921504606846976 is 2 to the 60~
     << cos(interval(1152921504606846976.0)) << endl;
@outputs^[-1, 0.5403023058681398]~
@outputs^[0.8414709848078965, 1]~
@outputs^[-inf, inf]~
@outputs^[-1.633123935319537e+16, 1.633123935319537e+16]~
@outputs^[-0.5567960822766418, -0.5567960822766416]~
\end{example}
\end{deffun}

\begin{deffun}{interval}{acos}{(const interval\& \var{I})}
  Returns the arccosine of \var{I}, on the part of \var{I} in \itv{-1}{1}.
\end{deffun}

\begin{deffun}{interval}{asin}{(const interval\& \var{I})}
  Returns the arcsine of \var{I}, on the part of \var{I} in \itv{-1}{1}.
\end{deffun}

\begin{deffun}{interval}{atan}{(const interval\& \var{I})}
  Returns the arctangent of \var{I}.
\begin{example}
cout << acos(interval(1,3)) << ' ' << asin(interval(1)) << ' '
     << atan(interval(1));
@outputs^[0] [1.570796326794896, 1.570796326794897]~
@outputs^[0.7853981633974482, 0.7853981633974484]~
\end{example}
\end{deffun}

\begin{deffun}{interval}{atan2}{(const interval\& \var{Y}, const interval\& \var{X})}
  Returns the hull of the angles in $(-\pi,\pi]$ of the points of ordinate
  in \var{Y} and abscissa in \var{X}, the point $(0,0)$ left out, as the
  \ieee{atan2} of IEEE~1788-2015 does: \code{atan2([0],[0])} is empty. The
  result is $[-\pi,\pi]$ whenever the box holds points of the half-line
  $y=0$, $x<0$, where the angle is $\pi$, and points below it, where the
  angle is next to $-\pi$. GAOL~4 declared the function, but it raised
  \code{unavailable\_feature\_error}.\newinvfive
\begin{example}
interval x(1,GAOL_INFINITY);
cout << atan2(interval(0.0), interval(0.0)) << endl
     << atan2(interval(-1,1), interval(-1)) << endl
     << atan2(x, x) << endl;
@outputs^[empty]~
@outputs^[-3.141592653589794, 3.141592653589794]~
@outputs^[0, 1.570796326794897]~
\end{example}
\end{deffun}

\deffunx{interval}{sinpi}{(const interval\& \var{I})}
\deffunx{interval}{cospi}{(const interval\& \var{I})}
\begin{deffun}{interval}{tanpi}{(const interval\& \var{I})}
  $\sin(\pi x)$, $\cos(\pi x)$ and $\tan(\pi x)$ for $x$ in \var{I}:
  \ieee{sinPi}, \ieee{cosPi} and \ieee{tanPi} of IEEE 1788-2015
  (Table~10.5).\newinvfive{} Their extrema and poles are at the multiples of
  $1/2$, which are doubles, and are found exactly, at every magnitude.
  \code{tanpi} has no value at the half-integers, which are left out: a
  half-integer alone gives the empty set.
\end{deffun}

\deffunx{interval}{asinpi}{(const interval\& \var{I})}
\deffunx{interval}{acospi}{(const interval\& \var{I})}
\deffunx{interval}{atanpi}{(const interval\& \var{I})}
\begin{deffun}{interval}{atan2pi}{(const interval\& \var{Y}, const interval\& \var{X})}
  $\operatorname{asin}(x)/\pi$, $\operatorname{acos}(x)/\pi$,
  $\operatorname{atan}(x)/\pi$ and $\operatorname{atan2}(y,x)/\pi$:
  \ieee{asinPi}, \ieee{acosPi}, \ieee{atanPi} and \ieee{atan2Pi} of IEEE
  1788-2015 (Table~10.5), within \itv{-1/2}{1/2}, \itv{0}{1},
  \itv{-1/2}{1/2} and \itv{-1}{1}.\newinvfive
\begin{example}
cout << sinpi(interval(1e17)) << ' '
     << sin(interval::pi() * 1e17) << endl;
@outputs^[0] [-1, 1]~
cout << tanpi(interval(0.5)) << ' '
     << tanpi(interval(0.25,0.75)) << ' '
     << asinpi(interval(1)) << ' '
     << atan2pi(interval(1), interval(-1)) << endl;
@outputs^[empty] [-inf, inf] [0.5] [0.75]~
\end{example}
\end{deffun}

\subsection{Hyperbolic functions}\label{sec:arithmetic-hyperbolic}
\cindex{hyperbolic functions}%

\begin{deffun}{interval}{cosh}{(const interval\& \var{I})}
  Returns the hyperbolic cosine of \var{I}.
\end{deffun}

\begin{deffun}{interval}{acosh}{(const interval\& \var{I})}
  Returns the hyperbolic arccosine of \var{I}, on the part of \var{I} in
  \itv{1}{+\infty}.
\end{deffun}

\begin{deffun}{interval}{sinh}{(const interval\& \var{I})}
  Returns the hyperbolic sine of \var{I}.
\end{deffun}

\begin{deffun}{interval}{asinh}{(const interval\& \var{I})}
  Returns the hyperbolic arcsine of \var{I}.
\end{deffun}

\begin{deffun}{interval}{tanh}{(const interval\& \var{I})}
  Returns the hyperbolic tangent of \var{I}, within \itv{-1}{1}.
\end{deffun}

\begin{deffun}{interval}{atanh}{(const interval\& \var{I})}
  Returns the hyperbolic arctangent of \var{I}, defined on
  $\left]-1,1\right[$, $-\infty$ and $+\infty$ being its limits at $-1$ and
  $1$. The part of \var{I} outside is left out, and an interval with no point
  of it gives the empty set, as \code{atanh(\itv{1}{5})} and
  \code{atanh(\itv{-1}{-1})} do.
\begin{example}
cout << cosh(interval(-1,2)) << ' '
     << tanh(interval(20,30)) << endl
     << acosh(interval(0,1)) << ' '
     << atanh(interval(-2,2)) << ' '
     << atanh(interval(1,5)) << endl;
@outputs^[1, 3.762195691083632] [0.9999999999999998, 1]~
@outputs^[0] [-inf, inf] [empty]~
\end{example}
\end{deffun}

\begin{Rem}\label{rem:arithmetic-hyperbolic-coremath}
The six hyperbolic functions are those of CORE-MATH. GAOL~4 took them from
the mathematical library of the system, which is sometimes further than one
double from the exact value.\newinvfive
\end{Rem}

\subsection{Other functions}\label{sec:arithmetic-other}

\begin{deffun}{interval}{hypot}{(const interval\& \var{X}, const interval\& \var{Y})}
  $\sqrt{x^2+y^2}$ for $x$ in \var{X} and $y$ in \var{Y}: \ieee{hypot} of
  IEEE 1788-2015 (Table~10.5), computed at the points of the box nearest to
  the origin and farthest from it.\newinvfive
\end{deffun}

\begin{deffun}{interval}{rsqrt}{(const interval\& \var{I})}
  $1/\sqrt{x}$ for $x$ in \var{I}, defined on $\left]0,+\infty\right]$,
  $+\infty$ being its limit at 0: \ieee{rSqrt} of IEEE 1788-2015
  (Table~10.5).\newinvfive
\begin{example}
cout << hypot(interval(3), interval(4)) << ' '
     << hypot(interval(-3,1), interval(4,5)) << endl
     << rsqrt(interval(4,16)) << ' '
     << rsqrt(interval(0.0,4)) << endl;
@outputs^[5] [4, 5.830951894845301]~
@outputs^[0.25, 0.5] [0.5, inf]~
\end{example}
\end{deffun}

\begin{deffun}{interval}{compoundm1}{(const interval\& \var{X}, const interval\& \var{Y})}
  $(1+x)^y-1$ for $x$ in \var{X} and $y$ in \var{Y}: \ieee{compoundm1} of
  IEEE 1788-2015 (Table~10.5). It keeps the digits of a small $x$ that
  \code{pow(1+x, y) - 1} loses, $1+x$ being rounded and the difference
  cancelling what is left. It is defined for $x>-1$, and for $x=-1$ when
  $y>0$, where it is $-1$. The part of the box outside is left out, and $-1$
  and $+\infty$, the limits at $x=-1$, at $x=+\infty$ and for an infinite
  $y$, are bounds. CORE-MATH has no function for it: the value at a corner
  of the box is computed with integers, 64 bits in a fast phase that settles
  71 to 95\,\% of the pairs, 128 bits and a bound of the error in the others.
  It takes 0.1 microsecond for each of the one or two corners where the
  bounds are, where the fast phase settles it, and 0.6 to 0.9 elsewhere. The
  bounds are the tightest ones, the double itself where the value is one
  (2.375 for $x=1.5-1$ and $y=3$, but also $2^m-1$ for $|m|\le53$, and the
  quotient $-x/(1+x)$ of $y=-1$, divided exactly), but in a case out of
  $10^{17}$ or so, where a bound is a double beyond.\newinvfive
\begin{example}
cout << compoundm1(interval(0.5), interval(2)) << ' '
     << compoundm1(interval(1,3), interval(1,2)) << endl
     << compoundm1(interval(-1), interval(0,1)) << ' '
     << compoundm1(interval(-1,0), interval(-1)) << endl
     << compoundm1(interval(1e-20), interval(3)) << endl;
@outputs^[1.25] [1, 15]~
@outputs^[-1] [0, inf]~
@outputs^[2.999999999999999e-20, 3.000000000000001e-20]~
\end{example}
\end{deffun}

\deffunx{interval}{erf}{(const interval\& \var{I})}
\begin{deffun}{interval}{erfc}{(const interval\& \var{I})}
  The error function $\operatorname{erf}(x)=\frac{2}{\sqrt{\pi}}\int_0^x
  e^{-t^2}\,dt$, increasing from $-1$ to 1, and its complement
  $\operatorname{erfc}(x)=1-\operatorname{erf}(x)$, decreasing from 2 to 0,
  for $x$ in \var{I}: the \code{erf} and \code{erfc} of C, which CORE-MATH
  computes, with the tightest bounds, exact at 0 and at the infinities.
  $\operatorname{erfc}(x)$ is subnormal from $x=26.55$ on, and below the
  least subnormal from $x=27.22$ on, where the lower bound is 0. IEEE
  1788-2015 does not name them, and \code{gaol\_ieee1788} does not hold
  them: a call \code{erf(\var{x})} on an interval finds them all the same,
  by argument-dependent lookup.\newinvfive
\begin{example}
cout << erf(interval(0.0)) << ' ' << erf(interval(-1,1)) << ' '
     << erf(interval(6,7)) << endl
     << erfc(interval(0.0,1)) << ' ' << erfc(interval(28,30)) << endl;
@outputs^[0] [-0.8427007929497149, 0.8427007929497149] [0.9999999999999998, 1]~
@outputs^[0.1572992070502851, 1] [0, 4.940656458412466e-324]~
\end{example}
\end{deffun}

\section{Relational arithmetic}\label{sec:relational-arithmetic}
\cindex{reverse functions|see{relational arithmetic}}%
\cindex{relational arithmetic}%

The relational functions of GAOL are the \dfn{reverse functions} of
IEEE 1788-2015 (10.5.4, Table~10.1), with their arguments in another order:
each computes the hull of the $x$ of an interval \var{I} whose image by the
function is in \var{J}. Their correspondence is given at the end of this
section (Table~\ref{tab:arithmetic-reverse}), and the namespace
\code{gaol\_ieee1788} provides them under their names
(Chapter~\ref{chap:ieee1788}).

\defmethodx{interval}{interval}{operator\% }{(double \var{d}) const}
\defmethodx{interval}{interval}{operator\% }{(const interval\& \var{I}) const}
\begin{deffun}{interval}{operator\%}{(double \var{d}, const interval\& \var{I})}
  \begin{operation}
    \op{$\code{self}\,\%\,\var{I}=\hull{\{z\in \RSet\mid
      \exists x\in \code{self}, \exists y\in I\colon x=yz\}}$}
  \end{operation}
  Relational division: \code{div\_rel(self, \var{I}, \itv{-\infty}{+\infty})}.
\begin{example}
cout << interval(1,2) % interval(0,1) << ' '
     << interval(1,2) % interval(-1,1);
@outputs^[1, inf] [-inf, inf]~
\end{example}
\end{deffun}

\subsection{$(n+1)$-ary relational functions}\label{sec:arithmetic-nary}

Consider the relation $y=\cos x$ between two interval variables $x$ and $y$.
The equivalent form $x=\operatorname{acos}y$ cannot be computed with the
function $\operatorname{acos}$, whose result always lies in \itv{0}{\pi}: a
relational version of $\operatorname{acos}$ is needed. For any value of $y$
in \itv{-1}{1} there are infinitely many $x$ such that $y=\cos x$, so the
domain of $x$ has to be taken into account, and the relational arccosine is
a binary operator.

\begin{Def}\label{def:arithmetic-acos-rel}
The \dfn{relational arccosine} of $Y$ with respect to $X$ is
\begin{equation*}
\code{acos\_rel}(Y,X)=\hull{\{x\in X\mid\exists y\in Y\colon y=\cos x\}}.
\end{equation*}
\end{Def}

\noindent It is to be contrasted with the functional arccosine,
\begin{equation*}
\code{acos}(Y)=\hull{\{x\in\RSet\mid\exists y\in Y\colon
  x=\operatorname{acos} y\}}.
\end{equation*}
Figure~\ref{fig:relational-cos} shows the results of $\code{acos}(J)$ and of
$\code{acos\_rel}(J,I)$ on the same arguments.

\begin{figure}[htbp]
  \centering
  \includegraphics[width=.85\textwidth,bb=11 262 553 470]{relation-cos}
  \caption{The functional arccosine $\code{acos}(J)$ and the relational
    arccosine $\code{acos\_rel}(J, I)$}
  \label{fig:relational-cos}
\end{figure}

\begin{Rem}\label{rem:arithmetic-empty-preimage}
An interval \var{I} that meets no piece of the preimage gives the empty
set.\newinvfive{} The same holds for every relational function of this
section: the result is empty when no $x$ of the interval contracted is in
the relation.
\end{Rem}

\begin{deffun}{interval}{acos\_rel}{(const interval\& J, const interval\& I)}
  \begin{operation}
    \op{$\code{acos\_rel}(J,I)=
      \hull{\{x\in I\mid\exists y\in J\colon y=\cos x\}}$}
  \end{operation}
  Returns the \emph{relational arccosine} of \var{J} with respect to \var{I}.
\end{deffun}

\begin{deffun}{interval}{asin\_rel}{(const interval\& J, const interval\& I)}
  \begin{operation}
    \op{$\code{asin\_rel}(J,I)=
      \hull{\{x\in I\mid\exists y\in J\colon y=\sin x\}}$}
  \end{operation}
  Returns the \emph{relational arcsine} of \var{J} with respect to \var{I}.
\end{deffun}

\begin{deffun}{interval}{atan\_rel}{(const interval\& J, const interval\& I)}
  \begin{operation}
    \op{$\code{atan\_rel}(J,I)=
      \hull{\{x\in I\mid\exists y\in J\colon y=\tan x\}}$}
  \end{operation}
  Returns the \emph{relational arctangent} of \var{J} with respect to
  \var{I}.

  The three functions enclose the multiples of $\pi$ of the pieces of the
  preimage in double-double arithmetic, and add the inverse function of
  \var{J} before rounding the sum: their bounds are within one double of the
  accurate ones.\newinvfive
\begin{example}
cout << asin_rel(interval(0.0), interval(1,7)) << endl
     << atan_rel(interval(1), interval(0,5)) << endl;
@outputs^[3.141592653589793, 6.283185307179588]~
@outputs^[0.7853981633974482, 3.926990816987242]~
\end{example}
\end{deffun}

\begin{deffun}{interval}{acosh\_rel}{(const interval\& J, const interval\& I)}
  \begin{operation}
    \op{$\code{acosh\_rel}(J,I)=
      \hull{\{x\in I\mid\exists y\in J\colon y=\cosh x\}}$}
  \end{operation}
  Returns the \emph{relational hyperbolic arccosine} of \var{J} with respect
  to \var{I}.
\end{deffun}

\deffunx{interval}{asinh\_rel}{(const interval\& J, const interval\& I)}
\begin{deffun}{interval}{atanh\_rel}{(const interval\& J, const interval\& I)}
  \begin{operation}
    \op{$\code{asinh\_rel}(J,I)=
      \hull{\{x\in I\mid\exists y\in J\colon y=\sinh x\}}$}
    \op{$\code{atanh\_rel}(J,I)=
      \hull{\{x\in I\mid\exists y\in J\colon y=\tanh x\}}$}
  \end{operation}
  Return the \emph{relational hyperbolic arcsine} and \emph{arctangent} of
  \var{J} with respect to \var{I}, that is, $\operatorname{asinh}(J)\cap I$
  and $\operatorname{atanh}(J)\cap I$, $\sinh$ and $\tanh$ being one-to-one.
  As $\tanh$ takes its values in $\left]-1,1\right[$, \code{atanh\_rel}
  gives the empty set for a \var{J} with no point of it, as \itv{1}{1}.
\begin{example}
cout << acosh_rel(interval(1,2), interval::universe()) << endl
     << asinh_rel(interval(0,1), interval::universe()) << endl
     << atanh_rel(interval(1), interval::universe()) << endl;
@outputs^[-1.316957896924817, 1.316957896924817]~
@outputs^[0, 0.8813735870195431]~
@outputs^[empty]~
\end{example}
\end{deffun}

\begin{deffun}{interval}{sqrt\_rel}{(const interval\& J, const interval\& I)}
  \begin{operation}
    \op{$\code{sqrt\_rel}(J,I)=\hull{\{x\in I\mid\exists y\in J\colon y=x^2\}}$}
  \end{operation}
  Returns the \emph{relational square root} of \var{J} with respect to
  \var{I} (Definition~\ref{def:arithmetic-relational-sqrt}).
\begin{example}
cout << sqrt_rel(interval(1,4), interval::universe()) << ' '
     << sqrt_rel(interval(1,4), interval(0,10));
@outputs^[-2, 2] [1, 2]~
\end{example}
\end{deffun}

\deffunx{interval}{nth\_root\_rel}{(const interval\& \var{J}, unsigned int \var{n}, \\\hbox{}\hfill const interval\& \var{I})}
\begin{deffun}{interval}{nth\_root\_rel}{(const interval\& \var{J}, int \var{p}, \\\hbox{}\hfill const interval\& \var{I})}
  \begin{operation}
    \op{$\code{nth\_root\_rel}(J,n,I)=
      \hull{\{x\in I\mid\exists y\in J\colon y=x^n\}}$}
    \op{$\code{nth\_root\_rel}(J,p,I)=\hull{\{x\in I\mid x^p\text{ is defined
      and in }J\}}$}
  \end{operation}
  Returns the \emph{relational inverse \var{n}-th power} of \var{J} with
  respect to \var{I}, for $n\geq1$. It takes the roots of
  \code{nth\_root()}.\newinvfive{} For an \code{int} \var{p}, or an integer
  of another type, it is the reverse of $x^p$ for every integer $p$, the
  \ieee{pownRev} of the standard. $x^0$ is $1$ for every $x$, $0$ included,
  so that the result is \var{I} where $1\in J$, and the empty set elsewhere:
  \code{nth\_root\_rel(\itv{-1}{5}, 0, I)} is \var{I}, and
  \code{nth\_root\_rel(\itv{2}{3}, 0, I)} is empty. For $p<0$, $x^p$ is
  $1/x^{|p|}$, which $0$ has not. \var{J} is then cut at $0$, and each part
  gives its roots: the hull of the roots of the whole of \var{J} would keep
  the part of \var{I} between them. The roots reach $0$ where \var{J} is
  unbounded, $0$ being a limit there and no solution:
  \code{nth\_root\_rel(\itv{1}{+\infty}, -2, \itv{0}{0})} is empty. The
  bounds are those of \code{nth\_root(J, p)}, the root and then its inverse:
  up to four doubles beyond the tightest ones for $p<0$. An exponent of type
  \code{double} or \code{float} is refused at compile time, as the ambiguous
  \code{nth\_root(x, 2.0)} is: it was converted to an \code{unsigned int},
  which cut $2.5$ to $2$.
\begin{example}
interval J(1,16);
cout << nth_root_rel(J, 4, interval::universe()) << ' '
     << nth_root_rel(J, 4, interval(-10,0.0));
@outputs^[-2, 2] [-2, -1]~
\end{example}
\begin{example}
cout << nth_root_rel(interval(1,8), -3, interval::universe()) << ' '
     << nth_root_rel(interval(-1,2), -3, interval(-0.5,0.5)) << ' '
     << nth_root_rel(interval(1,GAOL_INFINITY), -2, interval(0,3)) << ' '
     << nth_root_rel(interval(1,2), 0, interval(-5,5));
@outputs^[0.5, 1] [empty] [0, 1] [-5, 5]~
\end{example}
\end{deffun}

\begin{deffun}{interval}{invabs\_rel}{(const interval\& J, const interval\& I)}
  \begin{operation}
    \op{$\code{invabs\_rel}(J,I)=\hull{\{x\in I\mid |x|\in J\}}$}
  \end{operation}
  Returns the \emph{relational inverse absolute value} of \var{J} with
  respect to \var{I}. The negative part of \var{J} holds no magnitude:
  GAOL~4 took it as magnitudes, and gave \itv{-2}{2} where the result is
  empty.\newinvfive
\begin{example}
cout << invabs_rel(interval(1,2), interval::universe()) << ' '
     << invabs_rel(interval(-2,-1), interval(-10,10));
@outputs^[-2, 2] [empty]~
\end{example}
\end{deffun}

\begin{deffun}{interval}{div\_rel}{(const interval\& K, const interval\& J, const interval\& I)}
  \begin{operation}
    \op{$\code{div\_rel}(K,J,I)=
      \hull{\{x\in I\mid\exists z\in K,\exists y\in J\colon z=xy\}}$}
  \end{operation}
  Returns the \emph{ternary relational division} of \var{K} by \var{J} with
  respect to \var{I}, which is empty when no $x$ of \var{I} is in the
  relation.\newinvfive
\begin{example}
cout << div_rel(interval(1,2), interval(-1,1),
                interval(0,10)) << ' '
     << div_rel(interval(5,6), interval(1,2),
                interval(10,20));
@outputs^[1, 10] [empty]~
\end{example}
\end{deffun}

\begin{deffun}{interval}{pow\_rel}{(const interval\& K, const interval\& J, const interval\& I)}
  \begin{operation}
    \op{$\code{pow\_rel}(K,J,I)=
      \hull{\{x\in I\mid\exists z\in K,\exists y\in J\colon z=x^y\}}$}
  \end{operation}
  Returns the \emph{relational base} of the power $z=x^y$: the hull of the
  $x$ of \var{I} with an $x^y$ in \var{K} for a $y$ of \var{J}. It contracts
  the base of $z=x^y$ as \code{div\_rel(K, J, I)} contracts a factor of
  $z=xy$, with the arguments in the same order.\newinvfive{} $x^y$ is the
  \code{pow} of IEEE 1788-2015 (Table~9.1), that of
  \code{gaol\_ieee1788::pow}: $e^{y\log x}$ for $x>0$, $0$ for $x=0$ and
  $y>0$, and no value elsewhere, so that no $x<0$ is a solution, in either
  namespace. The \code{pow} of \code{gaol} takes the integer power for a
  degenerate integer exponent $[n]$, which has negative bases too
  (Remark~\ref{rem:arithmetic-pow-monotone}): the reverse of $x^n$ is
  \code{nth\_root\_rel(K, n, I)}.
\begin{example}
cout << pow_rel(interval(4,9), interval(2), interval::universe()) << ' '
     << pow_rel(interval(2,3), interval(1,GAOL_INFINITY), interval(0,1)) << ' '
     << pow_rel(interval(1,4), interval(2), interval(-3,3));
@outputs^[2, 3] [empty] [1, 2]~
\end{example}
  $x^y\geq2$ for a $y\geq1$ asks $x>1$: the second result is empty.
\end{deffun}

\begin{deffun}{interval}{pow\_exponent\_rel}{(const interval\& K, const interval\& H, \\\hbox{}\hfill const interval\& I)}
  \begin{operation}
    \op{$\code{pow\_exponent\_rel}(K,H,I)=\hull{\{y\in I\mid\exists z\in K,
      \exists x\in H\colon z=x^y\}}$}
  \end{operation}
  Returns the \emph{relational exponent} of the power $z=x^y$, the same
  \code{pow} as \code{pow\_rel()}: the hull of the $y$ of \var{I} with an
  $x^y$ in \var{K} for an $x$ of \var{H}.\newinvfive
\begin{example}
cout << pow_exponent_rel(interval(8), interval(2), interval::universe()) << ' '
     << pow_exponent_rel(interval(3), interval(2,4), interval::universe());
@outputs^[3] [0.792481250360578, 1.584962500721157]~
\end{example}
  The bounds of both functions are the tightest ones, but in one case, given below. The ends of the
  solutions, $c^{1/y}$ and $\log c/\log x$ for bounds $c$, $x$ and $y$ of
  the arguments, are bounded by the largest and the smallest doubles $d$
  whose powers $d^y$, or $x^d$, the \code{pow} of CORE-MATH, correctly
  rounded, proves to be on either side of $c$. A bound of \var{I} where the
  solutions start between it and the next double is kept when its power is
  within one double of $c$, which only knowing whether that power is a
  double would settle:
  \code{pow\_rel(\itv{1+2^{-52}}{1+2^{-52}}, \itv{0.5}{0.5},
  \itv{0}{1+2^{-51}})} is \itv{1+2^{-51}}{1+2^{-51}}, where
  $(1+2^{-52})^2=1+2^{-51}+2^{-104}$ and no $x$ of \itv{0}{1+2^{-51}} is a
  solution: the tightest result for \var{I} one double wider.
\end{deffun}

\begin{Ex}\label{ex:arithmetic-pow-rel}
The two reverse functions of the power contract in turn the base and the
exponent of the relation $z=x^y$, each from the current domains of the two
others.
\begin{example}
interval x(0, 10), y(2, 3), z(4, 9);
x = pow_rel(z, y, x);            @rem^// the x of x with pow(x, y) in z for a y of y~
y = pow_exponent_rel(z, x, y);   @rem^// the y of y with pow(x, y) in z for an x of x~
cout << x << ' ' << y << endl;   @rem^// [1.587401051968199, 3] [2, 3]~
gaol::cleanup();
\end{example}
The base is contracted to \itv{1.587401051968199}{3}, the cube root of 4
being its least value with a power in \itv{4}{9}; the exponent \itv{2}{3} is
not contracted, every $y$ of it having a base in \itv{1.587401051968199}{3}.
\end{Ex}

\deffunx{interval}{max\_rel}{(const interval\& Z, const interval\& Y, const interval\& X)}
\begin{deffun}{interval}{min\_rel}{(const interval\& Z, const interval\& Y, const interval\& X)}
  \begin{operation}
    \op{$\code{max\_rel}(Z,Y,X)=\hull{\{x\in X\mid\exists y\in Y\colon
      \max(x,y)\in Z\}}$}
    \op{$\code{min\_rel}(Z,Y,X)=\hull{\{x\in X\mid\exists y\in Y\colon
      \min(x,y)\in Z\}}$}
  \end{operation}
  Return the \emph{relational maximum} and \emph{minimum} of \var{Z} given
  \var{Y} with respect to \var{X}, the arguments in the order of
  \code{div\_rel}: the result, the other operand, then the operand
  contracted. $\max$ and $\min$ being symmetric, \code{max\_rel(Z,~X,~Y)}
  contracts \var{Y}, and \code{max(X,~Y)~\&~Z} contracts \var{Z}.
  $\max(x,y)$ is in \var{Z} when $x$ and $y$ are at most $\rightBound{Z}$
  and one of them at least $\leftBound{Z}$: the set is empty when \var{Y}
  lies above \var{Z}, $X\cap\itv{-\infty}{\rightBound{Z}}$ when \var{Y}
  meets \var{Z}, and $X\cap Z$ when \var{Y} lies below \var{Z}, the maximum
  being $x$; \code{min\_rel} is its mirror image. The bounds are those of
  \var{X} and \var{Z}: each result is the tightest enclosure.\newinvfive
\begin{example}
cout << max_rel(interval(1,2), interval(3,4), interval(-5,5)) << ' '
     << max_rel(interval(1,2), interval(0,1), interval(-5,5)) << ' '
     << max_rel(interval(1,2), interval(-1,0), interval(-5,5)) << endl
     << min_rel(interval(1,2), interval(2,3), interval(-5,5)) << endl;
@outputs^[empty] [-5, 2] [1, 2]~
@outputs^[1, 5]~
\end{example}
\end{deffun}

\begin{deffun}{interval}{sign\_rel}{(const interval\& Z, const interval\& X)}
  \begin{operation}
    \op{$\code{sign\_rel}(Z,X)=
      \hull{\{x\in X\mid \operatorname{sign}(x)\in Z\}}$}
  \end{operation}
  Returns the \emph{relational inverse sign} of \var{Z} with respect to
  \var{X}: the negative $x$ of \var{X} where \var{Z} holds $-1$, 0 where
  \var{Z} and \var{X} hold it, and the positive $x$ of \var{X} where \var{Z}
  holds 1; the empty set where \var{Z} holds none of them, as
  \itv{0.5}{0.7}. The negative $x$ go up to 0, which their tightest
  enclosure takes although $\operatorname{sign}(0)=0$, the negative doubles
  next to 0 being in the set.\newinvfive
\begin{example}
cout << sign_rel(interval(-1), interval(-2,3)) << ' '
     << sign_rel(interval(-1), interval(0,3)) << ' '
     << sign_rel(interval(0,1), interval(-2,3)) << endl;
@outputs^[-2, 0] [empty] [0, 3]~
\end{example}
\end{deffun}

\begin{deffun}{interval}{floor\_rel}{(const interval\& Z, const interval\& X)}
  \begin{operation}
    \op{$\code{floor\_rel}(Z,X)=\hull{\{x\in X\mid \lfloor x\rfloor\in Z\}}$}
  \end{operation}
  Returns the \emph{relational inverse floor} of \var{Z} with respect to
  \var{X}: the part of \var{X} in
  $\left[\lceil\leftBound{Z}\rceil,\lfloor\rightBound{Z}\rfloor+1\right[$,
  the $x$ whose floor is an integer of \var{Z}, and the empty set when
  \var{Z} holds no integer. That interval is open at its upper end, which
  its tightest closed enclosure takes all the same, the reals just below it
  being in the set: \code{floor\_rel(interval(0), interval(0.5,~2))} is
  \itv{0.5}{1}, and \code{floor\_rel(interval(0), interval(1,~2))} is empty,
  \var{X} starting at that end. Beyond $2^{53}$, where every double is an
  integer, $\lfloor\rightBound{Z}\rfloor+1$ is no double, and the bound is
  the double above it: \code{floor\_rel} of \itv{2^{53}}{2^{53}} ends at
  $2^{53}+2$, and at $+\infty$ for \itv{\text{DBL\_MAX}}{\text{DBL\_MAX}},
  where \var{X} goes on to it. Each result is the tightest
  enclosure.\newinvfive
\begin{example}
cout << floor_rel(interval(0,2), interval(-5,5)) << ' '
     << floor_rel(interval(0), interval(1,2)) << ' '
     << floor_rel(interval(0.5,0.7), interval(-5,5)) << endl;
@outputs^[0, 3] [empty] [empty]~
\end{example}
\end{deffun}

\deffunx{interval}{atan2\_rel\_y}{(const interval\& \var{C}, const interval\& \var{X}, \\\hbox{}\hfill const interval\& \var{Y})}
\begin{deffun}{interval}{atan2\_rel\_x}{(const interval\& \var{C}, const interval\& \var{Y}, \\\hbox{}\hfill const interval\& \var{X})}
  \begin{operation}
    \op{$\code{atan2\_rel\_y}(C,X,Y)=\hull{\{y\in Y\mid\exists x\in X\colon
      \operatorname{atan2}(y,x)\in C\}}$}
    \op{$\code{atan2\_rel\_x}(C,Y,X)=\hull{\{x\in X\mid\exists y\in Y\colon
      \operatorname{atan2}(y,x)\in C\}}$}
  \end{operation}
  Return the \emph{relational} \code{atan2} with respect to its first and to
  its second argument: the hulls of the ordinates and of the abscissas of
  the points of the box $\var{X}\times\var{Y}$ whose angle is in \var{C},
  the point $(0,0)$, which has no angle, left out.\newinvfive{} As in
  \code{div\_rel}, the interval contracted comes last, and the other
  argument of \code{atan2} before it. The angles are those of \code{atan2},
  in $\left]-\pi,\pi\right]$: the half-line $y=0$, $x<0$ has the angle
  $\pi$, and a \var{C} holding $-\pi$ but not $\pi$ leaves it out; the
  angles of \var{C} beyond $\pi$ are no angle of any point. The bounds are
  within two doubles of the tightest ones.
\begin{example}
interval c(1,2), x(-1,1), y(0,1);
cout << atan2_rel_y(c, x, y) << ' '
     << atan2_rel_x(c, y, x) << endl
     << atan2_rel_x(interval(-4,-3), interval(0.0), x)
     << endl;
@outputs^[0, 1] [-0.4576575543602859, 0.6420926159343309]~
@outputs^[empty]~
\end{example}
\end{deffun}

\subsection{The reverse functions of the standard}
\label{sec:arithmetic-reverse-names}

Table~\ref{tab:arithmetic-reverse} gives the correspondence between the
reverse functions of IEEE 1788-2015 and the relational functions of GAOL.
The last argument $x$ of a function of the standard is
\itv{-\infty}{+\infty} when it is left out.

\begin{table}[htbp]
  \centering\renewcommand{\arraystretch}{1.1}
  \caption{The reverse functions of IEEE 1788-2015 and the relational functions of GAOL} \label{tab:arithmetic-reverse}
  \begin{tabular}{ll}
    \hline
    \multicolumn{1}{c}{IEEE 1788-2015} & \multicolumn{1}{c}{GAOL}\\
    \hline\hline
    \ieee{sqrRev}$(c,x)$ & \code{sqrt\_rel(c, x)}\\
    \ieee{absRev}$(c,x)$ & \code{invabs\_rel(c, x)}\\
    \ieee{pownRev}$(c,x,p)$ & \code{nth\_root\_rel(c, p, x)}\\
    \ieee{sinRev}, \ieee{cosRev}, \ieee{tanRev}$(c,x)$ &
      \code{asin\_rel}, \code{acos\_rel}, \code{atan\_rel(c, x)}\\
    \ieee{coshRev}$(c,x)$ & \code{acosh\_rel(c, x)}\\
    \ieee{mulRev}$(b,c,x)$ & \code{div\_rel(c, b, x)}\\
    \ieee{powRev1}$(b,c,x)$ & \code{pow\_rel(c, b, x)}\newinvfive\\
    \ieee{powRev2}$(a,c,x)$ & \code{pow\_exponent\_rel(c, a, x)}\newinvfive\\
    \ieee{atan2Rev1}$(b,c,x)$ & \code{atan2\_rel\_y(c, b, x)}\newinvfive\\
    \ieee{atan2Rev2}$(a,c,x)$ & \code{atan2\_rel\_x(c, a, x)}\newinvfive\\
    \hline
  \end{tabular}

\end{table}

\begin{Rem}\label{rem:arithmetic-reverse-names}
\ieee{mulRevToPair}$(b,c)$, the two parts of \code{div\_rel(c, b)}, is in
the namespace \code{gaol\_ieee1788} under its name
(Chapter~\ref{chap:ieee1788}).\newinvfive{} \ieee{pownRev} takes every
integer $p$, $p\leq0$ included. \code{max\_rel}, \code{min\_rel},
\code{sign\_rel} and \code{floor\_rel} are the reverse functions of $\max$,
$\min$, $\operatorname{sign}$ and $\operatorname{floor}$ as the standard
defines them (10.5.4), which it neither names nor requires, Table~10.1 not
listing them: the namespace \code{gaol\_ieee1788} has no name for them. A
program that opens that namespace calls them as \code{max\_rel(z, y, x)}
all the same, argument-dependent lookup finding them in
\code{gaol\_core}.\newinvfive
\end{Rem}

\section{Cancellative addition and subtraction}
\label{sec:arithmetic-cancellative}
\cindex{cancellative operations}%

The two functions of this section are the \ieee{cancelMinus} and
\ieee{cancelPlus} of IEEE 1788-2015 (10.5.6): \code{cancel\_minus(X, Y)} is
the tightest interval that, added to \var{Y}, gives an interval containing
\var{X}, when such an interval exists.

\begin{deffun}{interval}{cancel\_minus}{(const interval\& \var{X}, const interval\& \var{Y})}
  The tightest interval $Z$ such that $Y+Z$ contains $X$,
  \itv{\leftBound{X}-\leftBound{Y}}{\rightBound{X}-\rightBound{Y}}, which
  exists when \var{X} and \var{Y} are bounded and \var{X} is at least as wide
  as \var{Y}; \itv{-\infty}{+\infty} otherwise, and the empty set for an
  empty \var{X} and a bounded \var{Y}: \ieee{cancelMinus} of IEEE 1788-2015
  (10.5.6). Whether \var{X} is at least as wide as \var{Y} is decided
  exactly, which rounding the differences of the bounds cannot
  do.\newinvfive
\end{deffun}

\begin{deffun}{interval}{cancel\_plus}{(const interval\& \var{X}, const interval\& \var{Y})}
  \code{cancel\_minus(\var{X}, -\var{Y})}: \ieee{cancelPlus} of IEEE
  1788-2015.\newinvfive
\begin{example}
cout << cancel_minus(interval(1,5), interval(0,2)) << ' '
     << cancel_minus(interval(1,2), interval(0,3)) << ' '
     << cancel_plus(interval(1,5), interval(0,2));
@outputs^[1, 3] [-inf, inf] [3, 5]~
\end{example}
\end{deffun}

\section{Discussion}\label{sec:arithmetic-discussion}

This chapter has given the reference of the operations on intervals. Every
functional operation computes the hull of the image of the part of its
argument within the domain of IEEE 1788-2015, and gives the empty set for an
argument with no point of that domain. The elementary functions take their
values from CORE-MATH, which is correctly rounded: their bounds are the
tightest ones, or within a double or two of them, where GAOL~4 depended on
the mathematical library of the system. The power is hybrid: an integer
exponent, whatever its type, takes the integer power \ieee{pown}, with the
negative bases, and any other exponent the real power \ieee{pow}, at the
price of a result that is not monotone for the inclusion of the exponent;
the powers of the two namespaces are compared in Chapter~\ref{chap:overview}.

\medskip
The relational functions are the reverse functions of the standard with the
arguments in the order of GAOL, the result first and the operand contracted
last. GAOL v5 adds the reverse functions of the power and of \code{atan2},
      which the standard names, and those of the maximum, the minimum, the sign and
      the floor, which it defines but does not name, and gives the empty set wherever no point of the operand
is in the relation.\newinvfive{} How tight each operation is, and the
conditions under which the bounds are the same on every machine, are the
subject of Chapter~\ref{chap:accuracy}; how the bounds are computed, of
Chapter~\ref{chap:computing}.

\medskip
Chapter~\ref{chap:functions} describes the functions on intervals that are
not arithmetic operations: the numeric and the integer functions of an
interval, its splitting, and the union and the intersection of intervals.
